% The De-margining Artifact
%
% Build:  make -C paper          (latexmk -pdf main.tex)
%
% Numbers are not typed into this file. Every quantity in the results section
% comes from generated/, written by scripts/make_tables.py out of
% results/cells.csv. Until the evaluation has been run those slots render as
% "[pending]", so the draft compiles at any stage without ever carrying a
% number that no run produced.
\documentclass[11pt,a4paper]{article}

\usepackage[T1]{fontenc}
\usepackage[utf8]{inputenc}
\usepackage[margin=1in]{geometry}
\usepackage{amsmath,amssymb}
\usepackage{booktabs}
\usepackage{makecell}
\usepackage{siunitx}
\usepackage{graphicx}
\usepackage{float}
\usepackage{caption}
\usepackage{microtype}
\usepackage{needspace}
\usepackage{etoolbox}
\usepackage[title,titletoc]{appendix}
\usepackage[numbers,sort&compress]{natbib}
\usepackage[hidelinks]{hyperref}
\usepackage{xcolor}
\usepackage{listings}
\usepackage{amsthm}
% Stops a table or figure from drifting past the section that refers to it.
\usepackage[section]{placeins}

\sisetup{detect-weight=true}

% Keep floats in the text stream. The default float-page threshold (0.5)
% is what produced the half-empty pages behind Table 6 and Figures 4--5.
\renewcommand{\topfraction}{0.90}
\renewcommand{\bottomfraction}{0.80}
\renewcommand{\textfraction}{0.07}
\renewcommand{\floatpagefraction}{0.80}
\setcounter{topnumber}{3}
\setcounter{bottomnumber}{2}
\setcounter{totalnumber}{5}
\setlength{\floatsep}{10pt plus 2pt minus 2pt}
\setlength{\textfloatsep}{12pt plus 2pt minus 4pt}
\setlength{\intextsep}{6pt plus 2pt minus 2pt}

\captionsetup{
  font=small,
  labelfont=bf,
  skip=6pt,
  justification=raggedright,
  singlelinecheck=false,
}

% Jump to the next page rather than leave a heading or a proposition
% stranded with one line of body, or a paragraph split so a single line
% sits alone overleaf.
\clubpenalty=10000
\widowpenalty=10000
\displaywidowpenalty=10000
\preto\section{\Needspace{8\baselineskip}}
\BeforeBeginEnvironment{lstlisting}{\Needspace{18\baselineskip}}

\newcommand{\inputtable}[1]{%
  {\small\setlength{\tabcolsep}{4.5pt}\input{#1}}%
}
% Scale a table down only when it would overflow the text block.
\newcommand{\inputwidetable}[1]{%
  \begingroup
  \small\setlength{\tabcolsep}{4pt}%
  \sbox{0}{\input{#1}}%
  \ifdim\wd0>\linewidth
    \resizebox{\linewidth}{!}{\usebox{0}}%
  \else
    \usebox{0}%
  \fi
  \endgroup
}

\lstset{
  basicstyle=\ttfamily\small,
  columns=fullflexible,
  keepspaces=true,
  breaklines=true,
  frame=single,
  framesep=3pt,
  xleftmargin=1em,
  xrightmargin=1em,
  language=Python,
  showstringspaces=false,
  commentstyle=\itshape\color{gray!70!black},
  keywordstyle=\bfseries,
}

\newtheorem{proposition}{Proposition}
\newtheorem{remark}{Remark}
\BeforeBeginEnvironment{proposition}{\Needspace{10\baselineskip}}

% Sections and subsections only, so the contents stay on one page.
\setcounter{tocdepth}{2}

\newcommand{\TBD}{\textcolor{red}{[pending]}}
\newcommand{\pendingtable}[1]{%
  \begin{center}\textcolor{red}{[table #1 pending: run \texttt{scripts/study.py evaluate},
  then \texttt{scripts/make\_tables.py}]}\end{center}}
\newcommand{\pendingfigure}[1]{%
  \begin{center}\textcolor{red}{[figure #1 pending: run \texttt{scripts/study.py evaluate},
  then \texttt{scripts/make\_figures.py}]}\end{center}}

% Quantities filled in from the run. Defaults keep the draft compiling.
\providecommand{\NFilesFetched}{525}
\providecommand{\NFilesRead}{514}
\providecommand{\NMatchesArchive}{176{,}629}
\providecommand{\NForecasts}{\TBD}
\providecommand{\NCells}{\TBD}
\providecommand{\NLeagues}{\TBD}
\providecommand{\NBooks}{\TBD}
\providecommand{\SpreadPinnacle}{\TBD}
\providecommand{\SpreadBetSixtyFive}{\TBD}
\providecommand{\SpreadOneXTwo}{\TBD}
\providecommand{\SpreadOverUnder}{\TBD}
\providecommand{\NSignFlips}{\TBD}
\providecommand{\PctSignFlips}{\TBD}
\providecommand{\SpearmanMarginSpread}{\TBD}
\providecommand{\SpearmanMarginSpreadLow}{\TBD}
\providecommand{\SpearmanMarginSpreadHigh}{\TBD}
\providecommand{\SpearmanBookSpread}{\TBD}
\providecommand{\SpearmanBookSpreadLow}{\TBD}
\providecommand{\SpearmanBookSpreadHigh}{\TBD}
\providecommand{\SpearmanClusterLow}{\TBD}
\providecommand{\SpearmanClusterHigh}{\TBD}
\providecommand{\NMarginCells}{\TBD}
\providecommand{\NMarginBooks}{\TBD}
\providecommand{\NPairedCells}{\TBD}
\providecommand{\ShareOneXTwoRoi}{\TBD}
\providecommand{\ShareOneXTwoRps}{\TBD}
\providecommand{\PairRoiDiff}{\TBD}
\providecommand{\PairRoiDiffLow}{\TBD}
\providecommand{\PairRoiDiffHigh}{\TBD}
\providecommand{\xifrozen}{\TBD}
\providecommand{\xihalflife}{\TBD}
\providecommand{\xirpsbest}{\TBD}
\providecommand{\xirpsspread}{\TBD}
\providecommand{\xipairedn}{\TBD}
\providecommand{\xirpsnodecay}{\TBD}
\providecommand{\xirpsfastest}{\TBD}
\IfFileExists{generated/summary.tex}{\input{generated/summary.tex}}{}
\IfFileExists{generated/calibration-macros.tex}{\input{generated/calibration-macros.tex}}{}

\title{The De-margining Artifact:\\
does a forecasting model's measured edge over the bookmaker\\
survive the choice of margin-removal transform?}
\author{Liam Abourousse}
\date{\today}

\begin{document}
\raggedbottom

% Title alone and vertically centred on the opening page. Written out rather
% than \maketitle: the article class sets the title block at the top of the
% page, and padding it with \vfill produces a stray page instead of centring it.
\begin{titlepage}
  \centering
  \vspace*{\fill}
  {\LARGE\bfseries The De-margining Artifact\par}
  \vspace{1.2em}
  {\large\itshape does a forecasting model's measured edge over the bookmaker\\
   survive the choice of margin-removal transform?\par}
  \vspace{3em}
  {\large Liam Abourousse\par}
  \vspace{0.8em}
  {\large\today\par}
  \vspace*{\fill}
\end{titlepage}
\setcounter{page}{2}

\begin{abstract}
Empirical football-forecasting studies routinely claim an edge over the
bookmaker. Substantiating such a claim requires converting offered odds into a
benchmark probability, which means removing the bookmaker's margin. At least
four transforms are in use for this: proportional normalisation, the power
transform, Shin's insider-trading model, and the constant odds-ratio
transform. Papers typically apply one without justification. The four are
not equivalent: they differ chiefly in how much margin they attribute to
longshots, which is precisely where a model's claimed edge tends to
concentrate.

We hold a single walk-forward Dixon--Coles model fixed across all four
transforms and vary nothing but the benchmark. The comparison covers
\NMatchesArchive{} archived matches, \NForecasts{} out-of-sample forecasts, and
\NCells{} (league, bookmaker, market) cells drawn from \NLeagues{} leagues and
\NBooks{} bookmakers, on 1X2 and Over/Under~2.5 closing prices. Three
predictions were pre-registered before the evaluation period was touched: that
transform disagreement scales with the book's margin (P1); that it is larger in
the three-outcome market than the two-outcome one (P2); and that there exist
cells in which the measured edge changes \emph{sign} between the proportional
and Shin benchmarks (P3). P1 is directionally supported: across the
\NMarginCells{} cells, Spearman's rank correlation between a cell's book margin
and the spread across transforms is \SpearmanMarginSpread{}. P2 is
directionally consistent with the prediction, although it rests on only
\NPairedCells{} paired cells. P3 does not hold, in any meaningful sense. Only
\NSignFlips{} cell of \NCells{} reverses sign, on
returns that are statistically indistinguishable from zero on both sides.
The fixed model achieves no edge anywhere in the sample, and loses to the
closing line in effectively every cell. The choice of transform demonstrably
changes what is measured. Whether it can overturn a published conclusion
remains open, and this design cannot settle it.
\end{abstract}
\clearpage

\tableofcontents
\clearpage

\section{Introduction}
\label{sec:intro}

A bookmaker's quoted prices do not sum to a probability. For a market with
outcomes $i = 1,\dots,n$ and offered decimal odds $d_i$, the raw implied
masses $p_i = 1/d_i$ sum to $B = \sum_i p_i > 1$; the excess $B - 1$ is the
overround, the bookmaker's margin. Any study that claims a forecasting edge
must first turn $\{p_i\}$ into a probability vector $\{\pi_i\}$ summing to one,
because it is against that vector, not against the quoted prices, that the
model's superiority is asserted.

That conversion is a modelling choice, and the literature has not converged on
one. Proportional normalisation is the default. Shin's model derives a
transform from the presence of insider money \citep{shin1991,shin1993}.
\citet{strumbelj2014} compared those two as probability estimators against two
logit regressions, and found Shin more accurate than basic normalisation for
every bookmaker/sport pair (Section~\ref{sec:related}). The power transform
\citep{clarke2017} and the constant odds-ratio transform \citep{cheung2015}
were not in that comparison. The four coincide only in the limit $B \to 1$. Away from it
they differ in how the removed mass is distributed across outcomes, and they
differ most at long odds. This matters because the favourite--longshot bias
concentrates market mispricing at exactly those prices \citep{cain2000}, and
because a model's largest apparent disagreements with the market are usually
found there too.

The question this paper asks is therefore not whether a model can beat the
market. It is narrower and, we argue, prior to it:

\begin{quote}
\emph{Does the measured edge of a fixed model over the bookmaker depend
materially on which de-margining transform constructs the benchmark?}
\end{quote}

If the answer is no, the field's casual treatment of the choice is harmless. If the answer is yes, and in particular if the sign of a measured edge can be
flipped by substituting one standard transform for another, then a class of
published results is under-determined by its own method section. Reporting
the transform then becomes a minimum requirement rather than a courtesy.

Our contribution is threefold. First, we isolate the transform as the sole
experimental factor: one model, one fit per matchday, one betting rule, four
benchmarks. Second, we pre-register the predictions and the frozen model
hyper-parameter before the evaluation period is run, so the git history of the
pre-registration timestamps the hypotheses ahead of the results
(Section~\ref{sec:prereg}). Third, we report the full coverage matrix,
including every excluded cell, so the reader can see what the study could not
measure as well as what it could.

\section{Related work}
\label{sec:related}

The question this paper isolates sits at the intersection of three literatures
that rarely cite one another: how a bookmaker's margin should be removed, how
a football match should be forecast, and how the resulting forecasts should be
compared against the market. Each is summarised in turn, and each summary ends
by saying what it does, and does not, settle for the design used here.

\subsection{Removing the margin, and where it goes}

Overround exists because the bookmaker prices risk, not merely outcomes,
and the earliest formal account of why is Shin's: modelling a market maker
who faces bettors who may be privately informed, Shin derives the odds a
risk-neutral bookmaker should set to protect against insider money and, by
inverting the relationship, a rule for recovering fair probabilities from
observed odds \citep{shin1991,shin1993}. The rule requires solving for a
single latent parameter, $z$, interpreted as the fraction of turnover
originating from insiders. The same $z$ later gets used as a measure of
market's informational efficiency in its own right, not merely as a
by-product of de-margining. Cortis derives the corresponding accounting
identity from the bookmaker's side, giving closed-form expressions for the
mean and variance of payout as a function of margin and stake distribution,
which is the complementary result to Shin's demand-side model: together they
describe the same over-round from the two sides of the book
\citep{cortis2015}.

Proportional normalisation has no comparable structural derivation. It makes
no assumption about how the margin is allocated across outcomes and preserves
the ratio of any two implied probabilities, which is why it is the field's
default. Which of these maps is the most accurate probability forecast of the
realised outcome has been studied at scale more than once.
\citet{strumbelj2014} treats the choice explicitly as a design decision and
compares basic (proportional) normalisation, Shin, and two logit regressions
on 37 competitions across five sports, finding Shin more accurate than
normalisation for every bookmaker/sport pair. \citet{goto2026} evaluate
multiplicative (proportional), Shin and power conversions on 90{,}014 football
matches across five bookmakers, and propose two further maps: an odds-only
equal-profitability method and a favourite--longshot generalised linear
model. Their question is which conversion performs best as a forecast. Ours
is whether, holding a forecasting model fixed, the scientific conclusion
about that model's edge over the bookmaker depends on the conversion. The two
contributions are complementary: a ranking of converters does not settle
whether swapping converters can change a model's measured edge, and a
sensitivity study does not settle which converter is the better probability
forecast. Neither paper asks whether the choice can reverse a model's
measured edge, which is the gap addressed here.

The reason any of this matters for a forecasting study is the
favourite--longshot bias: bettors accept a lower expected return on
longshots than on favourites, a regularity documented across parimutuel and
fixed-odds markets alike and surveyed at length by
\citet{ottaviani2008}, who catalogue explanations ranging from
risk-loving preferences to informational asymmetry of exactly the kind
Shin's model formalises. Whatever its cause, the bias means the market's
mispricing, and so a model's largest apparent disagreements with the
market, concentrates at long odds \citep{cain2000}. The four transforms
compared here differ most exactly there (Section~\ref{sec:transforms}), which
is the structural reason a transform choice can matter for a claimed edge
and not an incidental one.

\subsection{Forecasting the match}

The Poisson framework used here begins with \citet{maher1982}, who modelled
each team's goals scored and conceded as independent Poisson variables
governed by attack and defence parameters, recovered by maximum likelihood
from a season of scores. \citet{dixon1997} correct that model's chief
empirical failure, an under-prediction of low-scoring draws, with the
$\tau$ adjustment used unchanged in this study (Section~\ref{sec:method}),
and add exponential time-weighting so that a team's estimated strength
reflects its recent form rather than an unweighted season average. Both
extensions are structural to the model fit here. Neither addresses how the
resulting forecast should be benchmarked, which is outside the scope of the
original paper and the reason this study exists as a separate exercise.

Two further lines extend Dixon--Coles in directions this study does not take.
\citet{rue2000} replace the static per-season attack and defence parameters
with a state-space model updated match-by-match via Markov chain Monte Carlo,
letting strength drift continuously rather than being refit from a fixed
window; \citet{karlis2003} replace the independent-Poisson-plus-correction
device with a bivariate Poisson distribution that models the scoreline
dependence directly through a shared latent component, dispensing with the
four-cell patch that $\tau$ represents. Both are more sophisticated models of
the scoreline than the one used here, and both would very plausibly forecast
better. Adopting either would improve the \emph{level} of the model's
accuracy without changing the paper's design, which needs the forecasting
model held fixed rather than optimised: see Section~\ref{sec:method} for why,
and Section~\ref{sec:limitations} for what a stronger model would and would
not add to this study's conclusion.

\subsection{Is the market beatable, and how would you know}

\citet{levitt2004} shows that bookmakers behave less like risk-neutral
market-clearing intermediaries than like informed traders who set prices to
exploit predictable bettor bias while carrying inventory risk themselves.
That is one structural reason the closing line is a demanding
benchmark, not merely a busy one. Directly testing whether it can be beaten
by a statistical model, \citet{forrest2005} compare bookmaker-implied
forecasts against an independent statistical model for English football and
find the bookmakers' forecasts equal or superior on most measures, a result
the present study's Table~\ref{tab:dm} reproduces in direction, though not
because a similarly specified model must always lose. A market can be
soft in a segment a study happens not to cover. \citet{strumbelj2010} extend
the comparison across ten online bookmakers and six leagues and find
material differences in forecast quality \emph{between} bookmakers, which is
part of the motivation for reporting results per bookmaker here rather than
pooling books together (Table~\ref{tab:p1}). \citet{franck2011} give one
reason those differences persist: bookmakers price in sentimental demand as
well as expected outcome, shading odds toward popular teams, which is a
source of cross-sectional heterogeneity in the margin that this study's
design absorbs into the per-cell margin rather than modelling directly.

None of these papers asks the question this one does. They study whether the
market can be beaten, or why bookmakers price the way they do. This paper
holds the answer to the first question fixed (our model does not beat the
market anywhere in the sample: Section~\ref{sec:results}) and asks
whether the \emph{measurement} of an edge, beaten or not, is itself stable
under a choice every one of these papers had to make and rarely reports.

\subsection{Scoring the forecast}

\citet{brier1950} and \citet{epstein1969} give football forecasting its two
standard proper scores, and \citet{gneiting2007} give the general theory of
why propriety (the forecaster is not rewarded for reporting anything other
than a genuine belief) is the property worth insisting on. Whether the
football literature should prefer Epstein's distance-sensitive ranked
probability score over the ordering-blind Brier score is contested rather
than settled: \citet{constantinou2012} argue for RPS on the grounds that a
near-miss forecast (predicting a draw when the away side won by one) should
be scored better than a distant one, while \citet{wheatcroft2021} argues that
distance-sensitivity rewards nothing a forecast user actually needs and
recommends against it on those grounds. This study reports RPS as its primary
score, following the majority convention in the sport, but decomposes Brier
in Section~\ref{sec:metrics} precisely because the Murphy decomposition
\citep{murphy1973} used there is defined for Brier and has no RPS analogue;
nothing in the P1--P3 comparison depends on resolving the RPS-versus-Brier
question either way, since both are computed identically across all four
transforms and the comparison is paired throughout.

One measure conspicuously absent from the results is a Kelly-optimal stake
\citep{kelly1956}. Kelly staking sizes each bet by the model's own edge and
so entangles the measured return with the size of the (here, generally
absent) edge in a way flat staking does not. Section~\ref{sec:metrics}'s
one-unit rule is chosen so that the ROI comparison between transforms
isolates the benchmark, not the staking rule, at the cost of not reporting
the return a Kelly-sized bettor would actually have achieved.

\section{The instrument: four de-margining transforms}
\label{sec:transforms}

Each transform maps the raw implied masses $p \in (0,1)^n$, with book sum
$B = \sum_i p_i$, to a probability vector $\pi$ with $\sum_i \pi_i = 1$.

\paragraph{Proportional.} Remove margin uniformly in proportion to implied
mass:
\begin{equation}
  \pi_i = \frac{p_i}{B}.
\end{equation}
It fits no free parameter and preserves the ratio of any two implied masses.

\paragraph{Power.} Solve for the exponent $k$ satisfying
\begin{equation}
  \sum_i p_i^{\,k} = 1, \qquad \pi_i = p_i^{\,k}.
\end{equation}
Because every $p_i \in (0,1)$, the sum is strictly decreasing in $k$, equal to
$n$ at $k = 0$ and to $B$ at $k = 1$; a margined book therefore has its unique
root at $k > 1$. The transform takes proportionally more margin from longshots,
in the direction implied by the favourite--longshot bias
\citep{clarke2017}.

\paragraph{Shin.} Model a proportion $z$ of the money wagered as coming from
insiders, and recover \citep{shin1993}
\begin{equation}
  \pi_i = \frac{\sqrt{z^2 + 4(1-z)\,p_i^2 / B} - z}{2(1-z)},
  \label{eq:shin}
\end{equation}
with $z$ chosen so that $\sum_i \pi_i = 1$. Note that as $z \to 0$ this tends
to $p_i / \sqrt{B}$, not to $p_i / B$: Shin does not nest proportional
normalisation at its boundary.

\paragraph{Odds-ratio.} Impose a constant ratio $c$ between the offered and
true odds of every outcome, $p_i/(1-p_i) = c\,\pi_i/(1-\pi_i)$, which rearranges
to
\begin{equation}
  \pi_i = \frac{p_i}{c(1-p_i) + p_i},
\end{equation}
with $c$ solving $\sum_i \pi_i = 1$. At $c = 1$ the map is the identity and
$\pi_i$ is decreasing in $c$, so a margined book has its root at $c > 1$
\citep{cheung2015}.

\subsection{Existence and uniqueness of the fitted parameter}
\label{sec:existence}

Each of the three fitted transforms is presented above as ``solve for $k$
(or $z$, or $c$) such that the recovered vector sums to one,'' which
presupposes that a solution exists and that the solver converges to the
right one rather than an arbitrary one. That is worth establishing rather
than assuming, since a solver silently converging to a spurious root would
corrupt every downstream number without producing an error anywhere.

\begin{proposition}[Power]
For $p \in (0,1)^n$ with $B = \sum_i p_i > 1$, there is exactly one
$k^\ast > 1$ with $\sum_i p_i^{\,k^\ast} = 1$.
\end{proposition}
\begin{proof}
Let $f(k) = \sum_i p_i^{\,k}$. Each term $p_i^{\,k} = \exp(k \log p_i)$ has
derivative $p_i^{\,k}\log p_i$, strictly negative since $p_i \in (0,1)$
implies $\log p_i < 0$; hence $f$ is strictly decreasing and continuous on
$\mathbb{R}$. $f(1) = B > 1$ and $f(k) \to 0$ as $k \to \infty$ (each term is
bounded above by $\max_i p_i^{\,k} \to 0$), so by the intermediate value
theorem $f$ crosses $1$ exactly once, at some $k^\ast > 1$ since $f(1) > 1$
and $f$ is decreasing.
\end{proof}

\begin{proposition}[Odds-ratio]
For $p \in (0,1)^n$ with $B > 1$, there is exactly one $c^\ast > 1$ with
$\sum_i p_i / (c^\ast(1-p_i) + p_i) = 1$.
\end{proposition}
\begin{proof}
Write $g(c) = \sum_i p_i / (c(1-p_i) + p_i)$. Each term has derivative
$-p_i(1-p_i) / (c(1-p_i)+p_i)^2 < 0$ for $p_i \in (0,1)$ and $c > 0$, so $g$
is strictly decreasing on $(0,\infty)$. $g(1) = B > 1$, and $g(c) \to 0$ as
$c \to \infty$ termwise, so as with the power transform a unique root
$c^\ast > 1$ exists by the intermediate value theorem.
\end{proof}

\begin{proposition}[Shin]
For $p \in (0,1)^n$ with $B > 1$, there exists $z^\ast \in (0,1)$ solving
$\sum_i \pi_i(z^\ast) = 1$, where $\pi_i(z)$ is given by
Equation~\eqref{eq:shin}.
\end{proposition}
\begin{proof}
Write $h(z) = \sum_i \pi_i(z)$. At $z \to 0^+$, $\pi_i(z) \to p_i / \sqrt{B}$
termwise (Section~\ref{sec:transforms}), so $h(0^+) = B / \sqrt{B} =
\sqrt{B} > 1$. As $z \to 1^-$, write $s = 1-z$ and expand
$\sqrt{z^2 + 4s\,p_i^2/B} = \sqrt{(1-s)^2 + 4s\,p_i^2/B}$ to first order in
$s$: the term is $1 - s + 2s\,p_i^2/B + O(s^2)$, so
$\pi_i(z) = [(1-s+2s\,p_i^2/B) - (1-s)]/(2s) + O(s) \to p_i^2 / B$. Hence
$h(1^-) = \sum_i p_i^2 / B$. Since $0 < p_i < 1$ for every $i$, $p_i^2 <
p_i$ termwise, so $\sum_i p_i^2 < \sum_i p_i = B$ and $h(1^-) < 1$. $h$ is
continuous on $(0,1)$, so by the intermediate value theorem it crosses $1$
somewhere in $(0,1)$, giving existence of $z^\ast$. Monotonicity of $h$,
which gives uniqueness rather than merely existence, is established in
\citet{shin1993}. The $z \to 0^+$ boundary is exercised directly by the test
suite (\texttt{test\_shin\_limit\_as\_z\_vanishes\_is\_p\_over\_sqrt\_b} in
\texttt{tests/test\_demargin.py}), which exists because an earlier version of
this code used the wrong limit, $p_i/B$, a bug that summed to exactly one and
so produced no error anywhere downstream. The $z \to 1^-$ boundary above was
checked against the implementation directly for this paper rather than
asserted in the suite.
\end{proof}

The bisection driver used for all three (Section~\ref{sec:numerics}) does not
rely on these propositions to run. It will return \emph{some} value
whether or not a root exists in the bracket it is given, so the brackets
themselves are chosen wide enough that existence is never in doubt
($[1, \infty)$ enlarged geometrically for $k$ and $c$; $[0, 1-10^{-12}]$,
which brackets a sign change by the argument above, for $z$), and the
propositions are what justify treating the bisection's output as \emph{the}
answer rather than \emph{an} answer.

\subsection{A closed form for two-outcome markets}
\label{sec:twooutcome}

The Over/Under~2.5 market used throughout this study (Section~\ref{sec:data})
has $n = 2$, and at $n=2$ Shin's transform simplifies enough to solve without
a root-find at all. \citet{clarke2017} already record that Shin and additive
margin removal coincide for two-outcome markets. We record the closed form
and show why the fitted insider fraction $z$ drops out of the recovered
probabilities, which also explains part of why P2
(Section~\ref{sec:results}) finds smaller transform disagreement in the
two-outcome market than the three-outcome one.

\begin{proposition}
For $n = 2$, write $p = (p_1, p_2)$, $B = p_1 + p_2 > 1$, $d = p_1 - p_2$.
The Shin probabilities do not depend on the fitted parameter $z$ and equal
\begin{equation}
  \pi_1 = \frac{1+d}{2} = p_1 - \frac{B-1}{2}, \qquad
  \pi_2 = \frac{1-d}{2} = p_2 - \frac{B-1}{2}.
\end{equation}
\end{proposition}
\begin{proof}
Summing the Shin equation over $i=1,2$ and clearing the common denominator
$2(1-z)$ gives
$\sqrt{z^2+4(1-z)p_1^2/B} + \sqrt{z^2+4(1-z)p_2^2/B} = 2$: the two linear-in-$z$
terms that appear for general $n$ collapse to a constant precisely because
$n = 2$. Write $s = 1-z$, $u = \sqrt{z^2+4s\,p_1^2/B}$, $v =
\sqrt{z^2+4s\,p_2^2/B}$, so $u + v = 2$. Since $B = p_1+p_2$,
\[
  u^2 - v^2 = \frac{4s(p_1^2-p_2^2)}{B} = \frac{4s(p_1-p_2)(p_1+p_2)}{B} = 4sd,
\]
and $u - v = (u^2-v^2)/(u+v) = 2sd$. Solving the pair $u+v=2$, $u-v=2sd$
gives $u = 1+sd$ and $v = 1-sd$. Substituting into
$\pi_1 = (u-z)/(2s)$ with $z = 1-s$:
\[
  \pi_1 = \frac{(1+sd) - (1-s)}{2s} = \frac{s(1+d)}{2s} = \frac{1+d}{2},
\]
and symmetrically $\pi_2 = (1-d)/2$. Neither expression involves $s$ or $z$:
whichever $z^\ast \in (0,1)$ solves the summed equation (Proposition~3), it
produces the same pair of probabilities, so the value of $z^\ast$ itself is
irrelevant to the recovered vector at $n=2$. Rewriting $(1+d)/2$ as $p_1 -
(p_1+p_2-1)/2 = p_1 - (B-1)/2$ gives the stated closed form. Verified against
the bisection output on 6 spot-checked books to a relative error below
$10^{-12}$.
\end{proof}

At $n=2$, then, Shin is exactly \emph{additive} margin removal: split the
excess $B-1$ evenly and subtract it from each raw implied probability,
independently of the insider-trading story that motivates the general-$n$
model. Proportional, power and odds-ratio remain \emph{multiplicative} at
every $n$, including $n=2$ (Table~\ref{tab:p2} shows they still disagree with
each other and with Shin there). The transform that models the
favourite--longshot asymmetry most explicitly is, at exactly two outcomes,
the one transform that cannot express any asymmetry between the two
outcomes' margin at all: $\pi_1 - \pi_2 = p_1 - p_2$ is preserved exactly,
whatever the margin. That is a structural reason, not merely an empirical
one, for P2's finding that disagreement is smaller in the two-outcome market.

\subsection{Numerical solution: a common bisection driver}
\label{sec:numerics}

Three of the four require a root-find. To ensure that no transform is
advantaged by a looser solver, all four share a single bisection driver at a
common tolerance of $10^{-12}$ with a common iteration budget, and each
transform's recovered vector is renormalised identically afterwards. The
fitted free parameter ($k$, $z$, $c$; $B$ for proportional, which fits nothing)
is retained per market so that the effort each transform expended on a given
book can be reported. Appendix~\ref{app:code} walks through the driver's
implementation and the numerical safeguards around it.

\section{Data}
\label{sec:data}

Match results and bookmaker odds come from Football-Data.co.uk
\citep{footballdata}, which publishes per-league, per-season CSV files
containing full-time scores and the closing prices of a panel of
bookmakers, and, for recent seasons, the opening prices as well.

\NFilesFetched{} season files were fetched and \NFilesRead{} parsed, yielding
\NMatchesArchive{} matches. The eleven unreadable files are recorded rather
than silently dropped: seven declare an internal division code that disagrees
with their filename, three carry structurally corrupted rows, and one declares
the same bookmaker column twice. An unreadable league-season is missing data,
and missing data nobody counted is how a coverage table becomes a lie.

Two markets are used: the three-outcome match-odds market (1X2) and the
two-outcome Over/Under~2.5 goals market. \emph{Closing} prices only. Opening
prices are never used as a benchmark, because a benchmark drawn before the
market has absorbed team news measures something other than the market's final
opinion. A match survives into a cell only if its closing book is
\emph{complete} (every outcome priced), since a transform applied to a partial
book normalises over the wrong support and returns a benchmark that is wrong
rather than merely noisy.

The unit of analysis is the (league, bookmaker, market) cell. Cells with fewer
than 200 matched fixtures are excluded. The threshold was fixed in the
pre-registration rather than chosen after inspection. The full coverage matrix
is published as a result, including every excluded cell
(Table~\ref{tab:coverage}).

\begin{table}[ht]
  \centering
  \caption{Coverage. Complete closing books per bookmaker and market, and the
  cells that clear the 200-fixture threshold.}
  \label{tab:coverage}
  \IfFileExists{generated/tab-coverage.tex}
    {\inputtable{generated/tab-coverage.tex}}
    {\pendingtable{coverage}}
\end{table}

\section{Method}
\label{sec:method}

\subsection{Forecasting model}
\label{sec:model}

The model is Dixon--Coles \citep{dixon1997}, building on the independent-Poisson
scoring model of \citet{maher1982}: independent Poisson scoring rates with a
low-score dependence correction. For a match between home team $h$ and away
team $a$,
\begin{equation}
  \lambda = \alpha_h \beta_a \gamma, \qquad \mu = \alpha_a \beta_h,
\end{equation}
where $\alpha$ is attack strength, $\beta$ defence weakness, and $\gamma$ a
per-league home advantage. The joint mass of a scoreline $(x, y)$ is
\begin{equation}
  P(x, y) = \tau(x, y, \lambda, \mu, \rho)\,
            \frac{e^{-\lambda}\lambda^x}{x!}\,
            \frac{e^{-\mu}\mu^y}{y!},
  \label{eq:dc-joint}
\end{equation}
with $\tau$ the Dixon--Coles low-score correction, a piecewise adjustment
applied to the four scorelines in which both sides score at most one and left
at $1$ elsewhere:
\begin{equation}
  \tau(x,y,\lambda,\mu,\rho) =
  \begin{cases}
    1 - \lambda\mu\rho & (x,y) = (0,0) \\
    1 + \lambda\rho     & (x,y) = (0,1) \\
    1 + \mu\rho         & (x,y) = (1,0) \\
    1 - \rho            & (x,y) = (1,1) \\
    1                   & \text{otherwise.}
  \end{cases}
  \label{eq:tau}
\end{equation}
A single dependence parameter $\rho$ governs the correction, and $\rho < 0$
inflates $\tau(0,0)$ and $\tau(1,1)$, correcting independent Poisson's chief
empirical failure, an under-prediction of low-scoring draws. Not every $\rho$
is admissible at a given $(\lambda,\mu)$: $\tau$ must stay non-negative at all
four corrected cells, which requires
\begin{equation}
  \rho \in \left[\max\left(-\frac{1}{\lambda}, -\frac{1}{\mu}\right),\
                  \min\left(\frac{1}{\lambda\mu}, 1\right)\right].
  \label{eq:rho-bounds}
\end{equation}
This region depends on the fixture's own rates, which is why a $\rho$ valid
for every match in a fit window can still fall outside it for a pairing not
in that window (Section~\ref{sec:walkforward}).

Fitting a matchday's model is weighted maximum likelihood over the log-scale
parameter vector $\theta = (\log\alpha, \log\beta, \log\gamma, \rho)$, with
match $m$ weighted by $w_m = \exp(-\xi\,t_m)$
(Section~\ref{sec:xi}) for its age $t_m$ in days at the fit cutoff. Dropping
the outcome-independent factorial terms, the objective is
\begin{equation}
  \ell(\theta) = \sum_m w_m \Big[
    \log\tau(x_m,y_m,\lambda_m,\mu_m,\rho)
    + x_m\log\lambda_m - \lambda_m
    + y_m\log\mu_m - \mu_m
  \Big],
  \label{eq:loglik}
\end{equation}
with $\lambda_m, \mu_m$ given by scaling each team's fitted attack and defence
parameters as above. The log parameterisation makes positivity structural
rather than a bound and gives $\partial\lambda_m/\partial\log\alpha_{h(m)} =
\lambda_m$ directly, which is what makes the gradient below closed-form rather
than merely differentiable in principle. Two further choices matter for
identifiability and are stated rather than left implicit: (i) scaling every
$\alpha$ by a constant $c$ and every $\beta$ by $1/c$ leaves every
$\lambda_m,\mu_m$ unchanged, so the likelihood has a flat direction. It is
removed by dropping one team's log-attack parameter from the free vector and
setting it to minus the sum of the rest, fixing $\sum_i \log\alpha_i = 0$;
(ii) the optimiser is started at $\rho_0 = -0.10$ rather than $0$, because the
$\tau(0,1)$ and $\tau(1,0)$ derivatives with respect to $\rho$ vanish
identically at $\rho = 0$ (Equation~\eqref{eq:tau-grad} below), so starting
exactly there gives the optimiser no first-order signal about $\rho$ from
those two cells.

Maximising Equation~\eqref{eq:loglik} needs its gradient, and needs it
several thousand times over the course of one walk-forward run
(Section~\ref{sec:walkforward}): one refit per matchday per league, roughly
40 matchdays a season across 33 seasons and 16 leagues. A finite-difference
gradient costs one likelihood evaluation per parameter, on the order of 80
for a 20-team league, once attack, defence, home advantage and $\rho$ are
counted, so an analytic gradient is not a micro-optimisation here but the
difference between a run measured in minutes and one measured in hours.
Differentiating Equation~\eqref{eq:tau} at the four corrected cells and zero
elsewhere gives
\begin{equation}
  \frac{\partial\tau}{\partial\lambda} =
  \begin{cases} -\mu\rho & (0,0) \\ \rho & (0,1) \\ 0 & \text{else,} \end{cases}
  \quad
  \frac{\partial\tau}{\partial\mu} =
  \begin{cases} -\lambda\rho & (0,0) \\ \rho & (1,0) \\ 0 & \text{else,} \end{cases}
  \quad
  \frac{\partial\tau}{\partial\rho} =
  \begin{cases}
    -\lambda\mu & (0,0) \\ \lambda & (0,1) \\ \mu & (1,0) \\ -1 & (1,1) \\
    0 & \text{else.}
  \end{cases}
  \label{eq:tau-grad}
\end{equation}
Writing the per-match contribution to $\partial\ell/\partial\log\lambda_m$ and
$\partial\ell/\partial\log\mu_m$ using $\partial\lambda_m/\partial\log\lambda_m
= \lambda_m$,
\begin{equation}
  s^\lambda_m = w_m\left[\frac{\lambda_m}{\tau_m}
    \frac{\partial\tau_m}{\partial\lambda_m} + x_m - \lambda_m\right],
  \qquad
  s^\mu_m = w_m\left[\frac{\mu_m}{\tau_m}
    \frac{\partial\tau_m}{\partial\mu_m} + y_m - \mu_m\right],
  \label{eq:score-terms}
\end{equation}
the gradient with respect to each team's log-attack and log-defence
parameter is a sum of $s^\lambda_m$ and $s^\mu_m$ over exactly the matches
that team played at home or away,
\begin{equation}
  \frac{\partial\ell}{\partial\log\alpha_i} =
    \sum_{m:\,h(m)=i} s^\lambda_m + \sum_{m:\,a(m)=i} s^\mu_m,
  \qquad
  \frac{\partial\ell}{\partial\log\beta_i} =
    \sum_{m:\,a(m)=i} s^\lambda_m + \sum_{m:\,h(m)=i} s^\mu_m,
  \label{eq:grad-team}
\end{equation}
with the two remaining parameters collecting a sum over every match rather
than a per-team subset,
\begin{equation}
  \frac{\partial\ell}{\partial\log\gamma} = \sum_m s^\lambda_m, \qquad
  \frac{\partial\ell}{\partial\rho} =
    \sum_m w_m\,\frac{\partial\tau_m/\partial\rho}{\tau_m}.
  \label{eq:grad-gamma-rho}
\end{equation}
Every sum in Equations~\eqref{eq:grad-team}
and~\eqref{eq:grad-gamma-rho} is a scatter-add keyed by team
index, computed without a Python-level loop over matches
(\texttt{numpy.bincount}; Appendix~\ref{app:code}). A naive implementation
looping over teams and matches is quadratic in the two. The scatter-add form
is a single linear pass over the matches instead. The gradient with respect
to the dropped attack parameter is
folded into the free entries by the chain rule implied by constraint~(i)
above: $\partial\ell/\partial(\text{free }\log\alpha_i) = \partial\ell/\partial
\log\alpha_i - \partial\ell/\partial\log\alpha_{\text{last}}$. The whole
gradient is verified against \texttt{scipy.optimize.check\_grad} in the test
suite on every commit, because a subtly wrong analytic gradient converges
silently to the wrong optimum: nothing downstream (not the likelihood
value, not the fitted probabilities, not any later test) would flag that
the optimum reached was wrong rather than merely different.

The model choice is deliberate and, for this study, incidental. A better model
(Rue and Salvesen's dynamically updated ratings, or Karlis and Ntzoufras's
bivariate Poisson, Section~\ref{sec:related}) would very plausibly be
better, and would change the level of every measured edge. The question here is
whether the \emph{measurement} moves when only the benchmark is swapped, and
that requires the model to be held fixed, not to be good.

\subsection{Walk-forward protocol}
\label{sec:walkforward}

Each league's seasons are partitioned in time into three disjoint periods,
fixed before any fitting:

\begin{center}
\begin{tabular}{lll}
  \toprule
  Period & Span & Use \\
  \midrule
  Burn-in     & first 3 seasons & fitting only; never forecast or scored \\
  Calibration & next 2 seasons  & used \emph{solely} to select $\xi$ \\
  Evaluation  & all remaining   & the study \\
  \bottomrule
\end{tabular}
\end{center}

Within the scored periods the protocol is: before each matchday, refit on every
match that kicked off \emph{strictly} earlier, weighted by
$w_m = \exp(-\xi\, t_m)$ where $t_m$ is the match's age in days at the cutoff;
predict that matchday, then roll forward. Burn-in matches remain in the fit window
throughout, excluded from scoring, not from estimation. A matchday whose fit
window holds fewer than 300 matches is skipped.

The strict inequality is the single point at which look-ahead leakage could
enter, so it is enforced in one place in the code and asserted directly by the
test suite.

A fixture is \emph{dropped rather than priced} in three circumstances: a team
absent from the fit window has no fitted strength (a promoted side); a fitted
rate falls outside $(0.05, 6.0)$ expected goals, which indicates an
unidentified parameter rather than a strong team; or $\rho$ is inadmissible for
that fixture's own $(\lambda, \mu)$. Admissibility depends on the fixture's own
rates, so a $\rho$ valid for every match in the fit window can still be invalid
for a pairing not in it, and the fit-time check is structurally unable to catch
that.

\subsection{The decay parameter, and how weakly it is identified}
\label{sec:xi}

$\xi$ was selected by minimising mean ranked probability score over the
calibration period only, then frozen at
\begin{equation}
  \xi = \xifrozen \quad (\text{half-life } \ln 2/\xi = 400 \text{ days}).
\end{equation}

Selection was \emph{paired}. Candidates do not price the same fixtures: a fit
can fail to converge, or leave a team unidentified, and an unpaired sweep
compared means over match sets differing by 43\% in size (7{,}804 to 11{,}152
forecasts), which is not a valid comparison: an arm that happens to drop harder
fixtures scores better for reasons unrelated to forecast quality. That sweep
placed its minimum at 700 days, on the edge of its grid. Scoring every
candidate on the \xipairedn{} fixtures that \emph{all} candidates forecast
moved the optimum to \xihalflife{} days and made it interior, with the
neighbouring candidates worse on either side (Table~\ref{tab:xi}).

\begin{table}[ht]
  \centering
  \caption{Paired selection of the decay rate on the calibration period.
  Every candidate is scored on the same \xipairedn{} fixtures.}
  \label{tab:xi}
  \IfFileExists{generated/tab-calibration.tex}
    {\inputtable{generated/tab-calibration.tex}}
    {\pendingtable{3}}
\end{table}

\begin{figure}[ht]
  \centering
  \IfFileExists{generated/fig-calibration.pdf}
    {\includegraphics[width=0.72\textwidth]{generated/fig-calibration.pdf}}
    {\pendingfigure{calibration}}
  \caption{Paired mean RPS against the decay half-life, log scale, for every
  candidate in the calibration sweep. The marked optimum sits at an interior
  half-life rather than at either edge of the grid, and the curve stays
  within a narrow band of the no-decay baseline over a wide neighbourhood
  around it: $\xi$ is chosen, but only weakly identified, exactly as
  Table~\ref{tab:xi} states in numbers.}
  \label{fig:calibration}
\end{figure}

Stated plainly: $\xi$ is weakly identified. Half-lives of 400 and 700 days
differ by $7.3 \times 10^{-5}$ in paired RPS, and anything between roughly 250
and 1200 days scores within $6 \times 10^{-4}$ of the optimum. We do not
present 400 days as a precisely determined quantity. What the curve does
establish is that the extremes are wrong. A 30-day half-life is far worse
(\xirpsfastest{}), and no decay at all is worse than moderate decay
(\xirpsnodecay{} against \xirpsbest{}), so exponential decay earns its place
without its exact rate
mattering much. Because one model is held fixed across all four transforms, a
weakly identified $\xi$ shifts every arm together and cannot generate a
difference between them.

\subsection{Metrics}
\label{sec:metrics}

\paragraph{Ranked probability score.} Forecast accuracy is scored with the
ranked probability score \citep{epstein1969}, which respects the ordering of
1X2 outcomes and is the appropriate proper score for this market
\citep{constantinou2012,gneiting2007}. For a forecast $\pi$ over $n$ ordered
categories and realised category $j$,
\begin{equation}
  \mathrm{RPS} = \frac{1}{n-1} \sum_{k=1}^{n-1}
    \left( \sum_{i \le k} \pi_i - \mathbf{1}[j \le k] \right)^2 .
\end{equation}
Both the model and each de-margined benchmark are scored on the same fixtures,
so the comparison is paired by construction.

\paragraph{Economic edge.} Prediction P3 is an economic statement, so it is
measured economically. The betting rule is deliberately minimal: back an
outcome when the model's probability exceeds the de-margined market probability
for that outcome, stake one unit, and settle at the \emph{offered} price. Only
the benchmark varies between arms. Varying the strategy as well would make any
difference in measured edge unattributable. Return on investment is total
profit over total turnover, and is reported as undefined, not zero, when a
transform selected no bets, since ``no bets placed'' and ``bets placed that
broke even'' are different findings.

Changing the transform therefore changes the \emph{set of bets}, not merely
the measurement of a fixed portfolio: an outcome that clears
$p_i^{\mathrm{model}} > p_i^{\mathrm{book},T}$ under proportional
normalisation may fail it under Shin, and conversely. That is the
economically relevant object --- it is how the benchmark affects an
actionable edge --- but it is not a transform-sensitivity of a held-fixed
portfolio. The complementary measure, which scores every fixture and
involves no selection, is the ranked probability score already reported
alongside ROI (Tables~\ref{tab:p1}--\ref{tab:dm}). P1 and P2 are stated on
both.

\paragraph{Score decomposition.} Where the source of a difference in accuracy
is at issue, the Brier score is decomposed following \citet{murphy1973}. We use
the exact five-term form
\begin{equation}
  \mathrm{BS} = \mathrm{REL} - \mathrm{RES} + \mathrm{UNC}
              + \mathrm{WBV} - 2\,\mathrm{COV},
\end{equation}
whose last two terms (within-bin forecast variance and within-bin
forecast--outcome covariance) are identically zero only for single-valued
bins. On binned continuous forecasts they are not; see
Appendix~\ref{app:murphy}.

\subsection{Inference}
\label{sec:inference}

Three tools, each chosen because the naive alternative is invalid on this data
rather than merely less powerful.

Forecast comparisons use the Diebold--Mariano test \citep{diebold1995} with
Newey--West HAC standard errors \citep{newey1987} at the standard truncation
lag $\lfloor 4(n/100)^{2/9} \rfloor$. Forecasts produced by overlapping fit
windows are serially correlated, and treating per-match score differences as
independent would understate the variance and overstate significance.

Interval estimates for ROI use a stationary block bootstrap
\citep{politis1994} with geometric blocks of mean length 20 and a fixed seed.
Betting returns are heavy-tailed and clustered by matchday. An i.i.d.\
bootstrap would resample away exactly the dependence that makes the true
interval wide.

The study runs a family of tests across leagues, bookmakers, markets and
transform pairs, so $p$-values are controlled for false discovery rate by the
Benjamini--Hochberg procedure \citep{benjamini1995} at $q = 0.10$ across that
family.

\section{Pre-registration}
\label{sec:prereg}

The predictions, the protocol, the exclusion threshold and the frozen $\xi$
were committed to the repository before the evaluation period was run. The git
timestamp on \texttt{paper/\allowbreak preregistration.md} is the claim. If that file's
history shows modification after the first evaluation run, every prediction
below should be treated as post-hoc. The evaluation script parses $\xi$ out of
the committed pre-registration and refuses to start unless it agrees with the
calibration output to $10^{-9}$. A rounded value in the document would be a
\emph{different} decay rate from the one selected, and registering one number
while running another would defeat the point of registering it.

\begin{center}
\begin{tabular}{p{0.06\textwidth}p{0.42\textwidth}p{0.42\textwidth}}
  \toprule
  ID & Prediction & Falsified by \\
  \midrule
  P1 & Disagreement between transforms scales with the book's margin:
       negligible on a $\approx$102\% book, material on a $\approx$107\% book &
       Transform spread on high-margin cells not exceeding that on low-margin
       cells \\
  P2 & Disagreement is larger for three-outcome 1X2 than for two-outcome
       Over/Under~2.5 &
       OU25 spread matching or exceeding 1X2 spread on the same cells \\
  P3 & There exist (league, bookmaker) cells where the measured edge
       \emph{changes sign} between proportional and Shin &
       No cell showing a sign reversal in ROI between those two transforms \\
  \bottomrule
\end{tabular}
\end{center}

P3 is the headline: sign reversal, not merely magnitude change, is what would
invalidate a published conclusion. We committed in advance to reporting all
three predictions with effect sizes and confidence intervals regardless of
outcome, to dropping or reframing none of them after seeing results, and to
writing a null on P3 with P1 and P2 holding as a complete paper.

\section{Results}
\label{sec:results}

The evaluation produced \NForecasts{} out-of-sample forecasts over \NCells{}
cells spanning \NLeagues{} leagues and \NBooks{} bookmakers.

\subsection{P1: disagreement scales with the book's margin}

P1 is directionally supported. Across the \NMarginCells{} cells, Spearman's rank
correlation between a cell's mean book sum and the spread of the four
transforms within that cell is \SpearmanMarginSpread{} (ordinary
pair-bootstrap 95\% interval [\SpearmanMarginSpreadLow{},
\SpearmanMarginSpreadHigh{}]). That interval treats cells as independent,
which they are not: the same bookmakers, and sometimes the same matches,
appear in multiple cells. A bookmaker-clustered bootstrap of the same
coefficient, which resamples books rather than cells, gives
[\SpearmanClusterLow{}, \SpearmanClusterHigh{}], which still excludes zero.
Aggregated to one point per
bookmaker --- the coarser and more conservative sample --- the relationship
is \SpearmanBookSpread{} (95\% interval [\SpearmanBookSpreadLow{},
\SpearmanBookSpreadHigh{}], $n = \NMarginBooks{}$), which is wide enough to
include zero (Figure~\ref{fig:margin-spread}). The pattern in Table~\ref{tab:p1} is monotone at the
extremes: the two sharpest books, Betfair Exchange at a mean sum of 1.008 and
Pinnacle at 1.031, both show a mean RPS spread of 0.0001, while books between
1.06 and 1.09 show 0.0004 to 0.0009. The transforms agree almost exactly where
there is little margin to remove and diverge by four to nine times as much
where there is more of it, which is what the limiting argument predicts.

The stated falsification condition was that the spread on Bet365 should fail to
exceed the spread on Pinnacle. It does exceed it, on both measures
(\SpreadBetSixtyFive{} against \SpreadPinnacle{} in ROI). The registered
contrast is therefore in the predicted direction; the evidence for a
precisely estimated correlation is not.

Two cautions. The ROI spread column is noisier than the RPS column and does not
order as cleanly, because ROI is a function of how many bets a cell generated
and the smallest cells here rest on three league-seasons. And the mechanism is
not in doubt (the transforms provably coincide as the book sum approaches
one), so what this measures is the size of the effect at realistic margins,
not whether it exists.

\begin{table}[H]
  \centering
  \caption{Transform disagreement against book margin. Spread is the range of
  a quantity across the four transforms within a cell.}
  \label{tab:p1}
  \IfFileExists{generated/tab-p1.tex}
    {\inputtable{generated/tab-p1.tex}}
    {\pendingtable{P1}}
\end{table}

\begin{figure}[ht]
  \centering
  \IfFileExists{generated/fig-margin-spread.pdf}
    {\includegraphics[width=0.62\textwidth]{generated/fig-margin-spread.pdf}}
    {\pendingfigure{margin-spread}}
  \caption{Mean bookmaker margin (book sum) against the mean ROI spread across
  the four transforms, one point per bookmaker, with a least-squares trend.
  $\rho$ and its 95\% bootstrap confidence interval are computed on these
  bookmaker-level points ($n = \NMarginBooks{}$). Because that is a coarser
  aggregation than the cell-level correlation quoted as
  \SpearmanMarginSpread{} above, the two
  numbers need not agree exactly, but both are positive and both put the
  higher-margin books toward the higher-spread end.}
  \label{fig:margin-spread}
\end{figure}

\begin{figure}[ht]
  \centering
  \IfFileExists{generated/fig-divergence.pdf}
    {\includegraphics[width=0.62\textwidth]{generated/fig-divergence.pdf}}
    {\pendingfigure{divergence}}
  \caption{The four transforms' recovered probabilities against the raw
  implied probability, for the single closing book with the highest measured
  margin in the archive, chosen by margin alone, before any transform was
  run, so the example is not cherry-picked for a dramatic-looking spread. The
  longshot outcome (smallest raw implied probability) is annotated with each
  transform's recovered value: this is where the transforms are expected, and
  shown, to disagree most, which is what no table of summary statistics can
  make visible.}
  \label{fig:divergence}
\end{figure}
\FloatBarrier

\subsection{P2: disagreement is larger in the three-outcome market}

P2 is directionally consistent with the prediction, on thin evidence. On the
cells where the same bookmaker prices both
markets, the mean ROI spread is \SpreadOneXTwo{} for 1X2 against
\SpreadOverUnder{} for Over/Under 2.5, and the 1X2 spread is the larger in
\ShareOneXTwoRoi{} of pairs on ROI and \ShareOneXTwoRps{} on RPS.

The evidence is thin and should be read as such. Only three bookmakers in the
archive carry closing prices for both markets, which leaves \NPairedCells{} paired
cells. The direction matches the prediction and the mechanism is clear: the
transforms differ chiefly in how they treat a longshot, and a near-even
two-outcome book does not have one. \NPairedCells{} pairs cannot distinguish a
moderate effect from a small one. A paired bootstrap of the mean ROI-spread
difference (1X2 minus Over/Under) is \PairRoiDiff{} (95\% interval
[\PairRoiDiffLow{}, \PairRoiDiffHigh{}]). That interval excludes zero, but it
treats the \NPairedCells{} (league, bookmaker) pairs as independent, which they
are not: only three bookmakers price both markets. It describes this sample;
it is not a licence to generalise.

\begin{table}[H]
  \centering
  \caption{Transform disagreement by market, on cells where both markets are
  available from the same bookmaker.}
  \label{tab:p2}
  \IfFileExists{generated/tab-p2.tex}
    {\inputtable{generated/tab-p2.tex}}
    {\pendingtable{P2}}
\end{table}

\subsection{P3: sign reversal of the measured edge}

\textbf{P3 is met on its letter and fails on its substance, and the failure is
the more informative of the two.}

Exactly \NSignFlips{} cell of \NCells{} (\PctSignFlips{}) reverses the sign of
measured ROI between proportional normalisation and Shin's method. On the
prediction as written (\emph{there exist} cells where the edge changes
sign), that is a positive result. Table~\ref{tab:p3} shows why it should not
be read as one.

The cell is B1/Coral/1X2 on 214 matches, barely above the 200-fixture
threshold fixed in the pre-registration. Proportional normalisation puts the
ROI at $+0.0020$ and Shin at $-0.0004$. Both bootstrap intervals span
approximately $\pm 0.10$. These are not two different conclusions about an
edge. They are the same conclusion (no measurable edge), with noise falling on
opposite sides of zero. A reversal between two quantities that are each
indistinguishable from zero carries no information about whether the choice of
transform can overturn a finding.

The deeper problem is that the premise of P3 does not hold in this data, for
reasons set out in the next section: the model has no edge anywhere, so there
is no edge whose sign could be reversed. P3 was constructed on the assumption
that a competently implemented Dixon--Coles model would show positive measured
edge in at least some cells, against at least the softer books. It does not.
The prediction was therefore not so much falsified as rendered untestable by
this design, which is a limitation of the study rather than a finding about
the transforms.

\begin{table}[H]
  \centering
  \caption{The one cell in which the sign of measured ROI differs across
  transforms.}
  \label{tab:p3}
  \IfFileExists{generated/tab-p3.tex}
    {\inputtable{generated/tab-p3.tex}}
    {\pendingtable{P3}}
\end{table}

\subsection{Forecast accuracy against the benchmark}

The model is beaten by the market comprehensively and almost uniformly.
Table~\ref{tab:dm} gives the picture: mean model RPS of 0.2053 against
benchmark RPS between 0.1987 and 0.1990, a gap of roughly 0.0065 in the
market's favour under every transform. The model achieves a lower RPS than the
de-margined benchmark in \emph{no} cell under power, Shin or the odds ratio,
and in one cell of \NCells{} under proportional normalisation. Around
three-quarters of cells reject equal predictive accuracy at $q<0.10$ after the
Benjamini--Hochberg correction, in every case in the market's favour.

\begin{table}[H]
  \centering
  \caption{Model versus benchmark RPS by transform, with Diebold--Mariano
  statistics and Benjamini--Hochberg adjusted $p$-values.}
  \label{tab:dm}
  \IfFileExists{generated/tab-dm.tex}
    {\inputwidetable{generated/tab-dm.tex}}
    {\pendingtable{DM}}
\end{table}

\begin{figure}[H]
  \centering
  \IfFileExists{generated/fig-cell-scatter.pdf}
    {\includegraphics[width=0.58\textwidth]{generated/fig-cell-scatter.pdf}}
    {\pendingfigure{cell-scatter}}
  \caption{Per-cell model RPS against de-margined market RPS, one point per
  (cell, transform) pair. Almost every point lies above the diagonal: the
  model loses to the market in the same pattern Table~\ref{tab:dm} reports
  in numbers.}
  \label{fig:cell-scatter}
\end{figure}

Economic returns agree. Mean ROI across cells ranges from $-7.9\%$ under
proportional normalisation to $-9.4\%$ under the power method, and between 9\%
and 11\% of cells show positive ROI at all. Exactly one cell of \NCells{} has a
bootstrap interval lying entirely above zero, which is fewer than the four or
five that would be expected by chance across this many cells.

\begin{figure}[H]
  \centering
  \IfFileExists{generated/fig-roi.pdf}
    {\includegraphics[width=\textwidth]{generated/fig-roi.pdf}}
    {\pendingfigure{roi}}
  \caption{Distribution of per-cell ROI, one panel per transform, with a
  zero line on each. The mass sits on both sides of zero under every
  transform. The model has no measurable edge against any of the four
  benchmarks, which is the finding that governs how P3 should be read.}
  \label{fig:roi}
\end{figure}

None of this is surprising and none of it indicates a defect. Dixon--Coles is
a 1997 model with two parameters per team, no player information, no market
data and no in-play adjustment, evaluated against closing prices that
aggregate everything the betting market learned before kickoff. That it loses
is the expected outcome. It is reported here because it determines what P3 can
mean, and because a study that measured an edge without checking whether one
existed would be reporting an artefact of its own arithmetic.

\section{Discussion}
\label{sec:discussion}

The methodological claim survives and the headline claim does not.

That the four transforms diverge, and that they diverge in proportion to the
margin being removed and more sharply where a longshot is present, is
established here on \NCells{} cells across \NLeagues{} leagues and
\NBooks{} bookmakers. A paper that de-margins a soft book by proportional
normalisation and one that uses Shin's method are not measuring the same
benchmark, and the gap between them is largest exactly where a model's claimed
edge tends to concentrate. Reporting which transform was used, and ideally
reporting a result under more than one, costs nothing and is not current
practice.

What this study cannot show is that the choice flips a conclusion, because our
model never reaches a conclusion worth flipping. The design held the model
fixed so that any difference between arms would be attributable to the
transform alone, which was the right control, but it assumed the fixed model
would land near the market. It does not come close. Against that backdrop the
single sign reversal we observe is noise around zero, and would be noise around
zero whichever transform produced it.

The honest reading is that this work establishes the mechanism and bounds its
size, while leaving the question it set out to answer open.

\subsection{An exploratory extension that would settle P3}
\label{sec:extension}

We do not rewrite the pre-registration, and we do not pretend a follow-up was
registered. P3 asked whether transform disagreement can reverse the sign of a
measured edge. The question remains open, because the model used here has no
edge to reverse. What would close the gap is an explicitly post-hoc extension
using either a stronger forecasting model (Rue and Salvesen's dynamically
updated ratings, or Karlis and Ntzoufras's bivariate Poisson,
Section~\ref{sec:related}), opening rather than closing lines, a deliberately
softer bookmaker or segment, or a published model that already claims a
positive edge. If that extension found a statistically credible edge that was
significant under transform $A$ and disappeared or reversed under transform
$B$, the headline claim would be settled. That is a different study; the
apparatus built here would run it unchanged.

\section{Limitations}
\label{sec:limitations}

The model is fixed by choice. This study says nothing about whether a better
model would beat the market. It asks only whether the measurement of any
model's edge depends on the transform. A model with a different error profile
would concentrate its disagreements with the market at different prices, and
the size, though not the existence, of the effect reported here is specific
to this model.

Only two markets and only closing prices are used. Nothing here generalises to
in-play markets, Asian handicaps, or opening lines.

Backtested returns are not achievable returns. There are no stake limits, no
line movement between the moment a price is observed and the moment a bet could
be struck, no commission, and no account restrictions. The ROI figures exist to
be compared \emph{across transforms}, which is what the design controls. Their
absolute level is not a claim about what a bettor would have earned.

The decay parameter is weakly identified (Section~\ref{sec:xi}). This does not
threaten the comparison, which holds it fixed across arms, but it does mean the
model should not be read as tuned.

The most consequential limitation is that the fixed model has no measurable
edge anywhere in the sample. P3 was framed on the assumption that it would show
positive edge in at least some cells against at least the softer books. It does
not, in any cell, under any transform. Everything reported about P3 is
therefore a statement about behaviour near zero rather than about whether the
transform choice can overturn a real finding. P1 and P2 are unaffected, since
both concern disagreement between benchmarks and neither requires the model to
be any good. The exploratory extension in Section~\ref{sec:extension} is the
natural next measurement; it is not part of this paper.

The calibration period is early. A fixed three-season burn-in and two-season
calibration on a 33-season archive places the selection of $\xi$ in 1996--97
for the longest-running leagues, up to 28 years before the evaluation period
ends. The split was fixed in the design specification before any data was
examined and was deliberately not revised once that consequence became
visible.

\section{Reproducibility}
\label{sec:repro}

All code, the pre-registration, and the build for this paper are in the
repository at \url{https://github.com/Bardakor/The-De-margining-Artifact}. The
pre-registration cites the repository's former name, which now redirects. That
file is left byte-for-byte as committed, because its value rests on not having
been edited after the fact. The study runs
in three stages (\texttt{coverage}, then \texttt{calibrate}, then
\texttt{evaluate}), in that order, with the pre-registration commit required between the second and
the third. Every stochastic step takes an explicit seed and no global random
state is touched. The tables in Section~\ref{sec:results} are generated
mechanically from the evaluation output rather than transcribed.

\clearpage
\thispagestyle{plain}
\vspace*{\fill}
\begin{center}
{\Huge\bfseries APPENDIX}
\end{center}
\vspace*{\fill}
\clearpage

\begin{appendices}

\section{The five-term Murphy decomposition}
\label{app:murphy}

Murphy's vector partition \citep{murphy1973} is usually quoted in three terms,
$\mathrm{BS} = \mathrm{REL} - \mathrm{RES} + \mathrm{UNC}$, which is exact only
when every bin holds a single distinct forecast value. Binning continuous
forecasts leaves a residual. Expanding $\sum_i (p_i - o_i)^2$ within a bin
around the bin means $\bar p, \bar o$ gives
\begin{equation}
  (p_i - o_i) = (p_i - \bar p) + (\bar p - \bar o) + (\bar o - o_i),
\end{equation}
and squaring leaves three squared terms plus one surviving cross term,
$-2(p_i - \bar p)(o_i - \bar o)$. Summed, that is minus twice the within-bin
\emph{covariance} between forecast and outcome. It vanishes only when, inside
every bin, higher forecasts carry no information about which events occurred.

The four-term form that stops at within-bin variance is therefore one term
short. It is easy to miss because worked examples with one or two points per
bin satisfy the zero-covariance condition by accident. On real continuous
forecasts it does not vanish: measured on 4{,}000 simulated forecasts the
covariance term accounted for 1.2\% of the Brier score, which the four-term
form silently misattributes to reliability and resolution. Including it drops
the reconstruction residual from order $10^{-3}$ to machine zero. All five
terms are reported, and the residual is computed explicitly so the identity can
be checked rather than trusted.

\section{Implementation notes: from equations to code}
\label{app:code}

The equations in Sections~\ref{sec:transforms} and~\ref{sec:model} describe
what is computed. This appendix shows how, at the four points where the
implementation choice is not a transcription of the mathematics but a
decision in its own right. Each excerpt is quoted verbatim from the
repository accompanying this paper (Section~\ref{sec:repro}) and is
shortened only by removing docstrings already quoted in the main text.

\subsection{A shared root-finder for three different equations}

Propositions 1--3 (Section~\ref{sec:existence}) establish that the power,
odds-ratio and Shin equations each have a unique root in a known bracket.
Rather than three solvers with three tolerances, one bisection driver is
shared, so that no transform's estimated parameter is more or less converged
than another's purely as an artefact of solver choice:

\begin{lstlisting}
def _bisect(objective: Callable[[float], float], lo: float, hi: float) -> float:
    low, high = lo, hi
    mid = 0.5 * (low + high)
    for _ in range(MAX_ITERATIONS):
        mid = 0.5 * (low + high)
        value = objective(mid)
        if abs(value) < TOLERANCE:
            return mid
        if value > 0.0:
            low = mid
        else:
            high = mid
    return mid
\end{lstlisting}
\emph{(\texttt{src/footy/market/demargin.py})}

\texttt{objective} is $f(k) - 1$, $g(c) - 1$, or $h(z) - 1$ from
Section~\ref{sec:existence}, each decreasing on its bracket by the
propositions proved there, so the sign of \texttt{value} always identifies
which half of the bracket contains the root. \texttt{TOLERANCE} is
$10^{-12}$ and \texttt{MAX\_ITERATIONS} is 200, both module constants shared
by every call site rather than parameters a caller could vary per transform.
Note what the function does \emph{not} do: it never raises when a root is
absent from the bracket, it simply returns whatever \texttt{mid} it last
tried. That is safe here only because Section~\ref{sec:existence} proves a
root always exists in the brackets actually passed in. \texttt{power} and
\texttt{odds\_ratio} grow their upper bound geometrically until the
objective's sign confirms it brackets the root before calling
\texttt{\_bisect} at all, and \texttt{shin} passes the fixed bracket
$[0, 1{-}10^{-12}]$ justified by Proposition~3's boundary values.

\subsection{Vectorising the Dixon--Coles gradient}

Equations~\eqref{eq:grad-team} sum $s^\lambda_m$ and $s^\mu_m$ over the
matches each team played at home or away. Written as a Python loop over
matches this is correct but, over several thousand refits, slow enough to
matter. Written as a scatter-add keyed by team index it is one call per
sum and no per-match Python overhead:

\begin{lstlisting}
grad_attack = np.bincount(data.home_idx, s_lam, minlength=n) + np.bincount(
    data.away_idx, s_mu, minlength=n
)
grad_defence = np.bincount(data.away_idx, s_lam, minlength=n) + np.bincount(
    data.home_idx, s_mu, minlength=n
)
grad_attack_free = grad_attack[:-1] - grad_attack[-1]
\end{lstlisting}
\emph{(\texttt{src/footy/fit/likelihood.py})}

\texttt{np.bincount(idx, weights, minlength=n)} returns, for each team index
$0,\dots,n-1$, the sum of \texttt{weights} at the positions where
\texttt{idx} equals that index: exactly $\sum_{m:\,h(m)=i} s^\lambda_m$
when \texttt{idx = home\_idx} and \texttt{weights = s\_lam}. The last line is
the chain-rule correction for the dropped attack parameter described after
Equation~\eqref{eq:grad-team}: every team's raw gradient is computed first,
and only then is the last team's gradient subtracted from the rest, because
the subtraction is a property of the \emph{parameterisation}, not of any
individual team, and computing it after the scatter-add keeps that fact
visible in the code rather than folded into per-team bookkeeping.

\subsection{The block bootstrap, vectorised over resamples}

Section~\ref{sec:inference} uses a stationary block bootstrap
\citep{politis1994} for ROI confidence intervals. Its defining mechanism
(continue the current block with probability $1 - 1/L$, or restart at a
uniformly random position with probability $1/L$, wrapping past the end of
the series) is inherently sequential in the resample \emph{index}, but not
across the $10{,}000$ resamples drawn at once, so the sequential part is a
Python loop over the (short) series length with every resample advanced
together as a vector:

\begin{lstlisting}
index = np.empty((n_resamples, n), dtype=np.int64)
index[:, 0] = starts[:, 0]
for t in range(1, n):
    advanced = (index[:, t - 1] + 1) % n
    index[:, t] = np.where(restart[:, t], starts[:, t], advanced)
\end{lstlisting}
\emph{(\texttt{src/footy/eval/inference.py})}

\texttt{restart} is a $(10{,}000, n)$ boolean array drawn once, true with
probability $1/L$ per entry, and \texttt{starts} a same-shaped array of
uniform random positions. The modulo in \texttt{advanced} is what makes a
block that runs past the last observation wrap around to the first, which is
the ``stationary'' in stationary bootstrap and what a naive block bootstrap
without wraparound would get wrong at the series boundary. The loop runs $n$
times, not $n \times 10{,}000$: every resample's index at step $t$ is
computed in one vectorised statement.

\subsection{Where the leakage check actually lives}

Section~\ref{sec:walkforward} states that a matchday's fit uses only matches
that kicked off \emph{strictly} before the cutoff, and that this is enforced
in one place rather than trusted at every call site:

\begin{lstlisting}
deltas = (pd.Timestamp(as_of) - stamps).dt.total_seconds().to_numpy(dtype=np.float64)
if np.any(deltas <= 0.0):
    n_bad = int((deltas <= 0.0).sum())
    raise ValueError(
        f"{n_bad} match(es) at or after the fit cutoff {as_of}; "
        "a fit may only use strictly earlier matches"
    )
\end{lstlisting}
\emph{(\texttt{src/footy/fit/decay.py}, inside the function that also
computes each match's age for the decay weight $w_m = \exp(-\xi t_m)$ of
Equation~\eqref{eq:loglik})}

Computing the decay weight and checking for leakage happen in the same
function rather than two separate ones because every caller that needs one
needs the other: any code path that asks ``how old is this match'' is, by
construction, also asked to answer ``could this match have leaked,'' and
putting both questions in one function makes it impossible to compute a
match's age without also having its non-negativity checked.

\end{appendices}

\clearpage

\bibliographystyle{plainnat}
\bibliography{refs}

\end{document}
